\subsection{Cartan Calculus}
In this section we discuss two differential operators on the graded algebra $\Omega^\bullet(M)$ of differential forms on a manifold $M$, namely the Lie derivative $\mathcal{L}_X$ and the interior product $\iota_X$, where $X$ is a smooth vector field on $M$. The main goal of this section is to show well-definedness of $\mathcal{L}_X$ and proof the following theorem, describing the relationships between $\mathcal{L}_X, \iota_X$ and the exterior derivative $d$.

\begin{theorem}
    Let $M$ be a manifold and $X,Y \in \mathscr{X}(M)$. Then the maps
    \begin{align*}
        d &\colon \Omega^r(M) \to \Omega^{r+1}(M), \\
        \mathcal{L}_X &\colon \Omega^r(M) \to \Omega^r(M), \\
        \iota_X &\colon \Omega^r(M) \to \Omega^{r-1}(M)
    \end{align*}
    are differential operators satisfying the following Leibniz rules for $\alpha, \beta \in \Omega^\bullet(M)$:
    \begin{align*}
        d(\alpha \wedge \beta) &= d\alpha \wedge \beta + (-1)^{|\alpha|} \alpha \wedge d\beta, \\
        \mathcal{L}_X(\alpha \wedge \beta) &= \mathcal{L}_X\alpha \wedge \beta +  \alpha \wedge \mathcal{L}_X\beta, \\
        \iota_X(\alpha \wedge \beta) &= \iota_X\alpha \wedge \beta + (-1)^{|\alpha|} \alpha \wedge \iota_X\beta,
    \end{align*}
    as well as the following identities:
    \begin{enumerate}[(i)]
        \item $d^2 = 0$,
        \item $d\mathcal{L}_X = \mathcal{L}_Xd$,
        \item $\mathcal{L}_X = d\iota_X + \iota_Xd$,
        \item $\mathcal{L}_X\mathcal{L}_Y - \mathcal{L}_Y\mathcal{L}_X = \mathcal{L}_{[X,Y]}$,
        \item $\mathcal{L}_X\iota_Y - \iota_Y\mathcal{L}_X = \iota_{[X,Y]}$,
        \item $\iota_X\iota_Y + \iota_Y\iota_X = 0$.
    \end{enumerate}
\end{theorem}

Firstly let us restate the definition of the Lie derivative. For a smooth boundaryless manifold $M$, $\alpha \in \Omega^r(M)$ and $X \in \mathscr{X}(M)$ we define
\[
(\mathcal{L}_X \alpha)_p := \frac{d}{dt}\bigg|_{t=0} (\phi_t^\ast \alpha)_p,
\]
where $\phi_t = \phi(t,-)$ denotes the flow of $X$.

\begin{lemma}
    \label{lem:well_definedness_Lie_derivative}
    $(\mathcal{L}_X\alpha)_p$ exists for all $p\in M$ and defines a smooth $r$-form on $M$.
\end{lemma}

Before showing the lemma, let us introduce the notion of a smooth family of $r$-forms.

\begin{definition}[Smooth family of $r$-forms]
    Let $\beta_t$ be a family of $r$-forms on $M$ indexed over some Interval $I$. We call $\beta_t$ smooth if in some (then all) local coordinates $x^1,\ldots,x^n$ we can write
    \[
    (\beta_t)_p = \sum \beta^{i_1,\ldots,i_r}(t,p)d_px^{i_1} \wedge \ldots \wedge d_px^{i_r}
    \]
    for smooth functions $\beta^{i_1,\ldots,i_r}$.
\end{definition}

\noindent Then $t \mapsto (\beta_t)_p$ is a smooth curve in the vector space $\mathrm{Alt}^r(T_pM)$, such that 
\[
\left(\frac{d}{dt}\bigg|_{t=0} \beta_t\right)_p := \frac{d}{dt}\bigg|_{t=0} (\beta_t)_p
\]
is a well-defined $r$-form. Furthermore, in local coordinates we have
\[
    \frac{d}{dt}\bigg|_{t=t_0} \beta_t = \sum \frac{\partial\beta^{i_1,\ldots,i_r}}{\partial t}(t_0,-)dx^{i_1} \wedge \ldots \wedge dx^{i_r},
\]
such that this $r$-form is alsow smooth.

\begin{proof}[Proof of Lemma \ref{lem:well_definedness_Lie_derivative}]
    
\end{proof}


\subsubsection{Moser`s trick}
